\bookchapter{第二十二章：AI如何在手機中工作}

\chapterintro{歡迎來到手機最酷的“魔術現場”，主角是：人工智慧。你每天在用的“一拍即得”的美照、張口就來的語音助手、越用越懂你的輸入法，背後都站著這位隱形大腦。}

\booksection{一、AI相機：你的“後期暗房”搬進了手機}

現在的AI計算攝影，更像是把一個專業修圖師+攝影師團隊直接塞進了快門鍵裡。在你按下快門前，AI就已經開始工作了：掃描場景，識別出“這是藍天”、“這是人臉”、“這是夜景”。

\booksection{二、語音助手：你的“聲音秘書”如何聽懂你？}

\begin{enumerate}
  \item \textbf{聽清（語音識別）}：將聲波轉變成文字。
  \item \textbf{聽懂（自然語言理解）}：理解意圖，提取關鍵實體。
  \item \textbf{執行與回答}：呼叫對應App或服務，再透過語音合成"說"給你聽。
\end{enumerate}

\booksection{三、端側AI vs. 雲端AI}

端側AI把計算留在手機裡，速度快、隱私高、可離線。雲端AI需要聯網，但能處理更復雜的任務。

\booksection{四、AI大模型與手機：你的隨身智慧助手}

2025-2026年，AI大模型（如DeepSeek R1、ChatGPT-5、文心一言5.0、通義千問3等）的爆發深刻改變了手機的使用方式。手機不再只是執行指令的工具，而是能理解、推理和創作的智慧夥伴。端側AI晶片（NPU算力突破80 TOPS）讓大模型推理首次在手機本地跑通，無需聯網也能用AI。

\booksubsection{什麼是AI大模型？}

AI大模型（Large Language Model，LLM）是透過海量文字資料訓練出來的AI系統，能理解和生成自然語言。你可以把它想象成一個讀過網際網路上幾乎所有文字的“超級學霸”，能回答問題、寫文章、翻譯、程式設計、分析資料。

\booksubsection{手機上的AI大模型應用：}

\begin{itemize}
  \item \textbf{對話助手}：DeepSeek、ChatGPT、豆包、Kimi等App，可以隨時提問、翻譯、寫文案、總結文章。比傳統搜尋引擎更智慧——它不只給你連結，而是直接給你答案。
  \item \textbf{寫作與創作}：讓AI幫你寫郵件、總結會議紀要、潤色文章、生成創意文案。
  \item \textbf{影象生成與編輯}：AI一鍵消除照片中的路人、擴充圖片邊界、風格轉換、甚至根據文字描述生成圖片。
  \item \textbf{實時翻譯}：通話實時翻譯、面對面對話翻譯、拍照翻譯，出國旅行不再有語言障礙。
  \item \textbf{學習輔導}：AI可以解釋概念、出題練習、批改作業，是隨身家教。
\end{itemize}

\booksubsection{端側大模型：手機自己跑AI}

最新的旗艦手機已經能在本地執行小型AI大模型（通常7B引數左右），帶來的變化是：

\begin{itemize}
  \item \textbf{不需要網路}：斷網也能用AI功能，飛機上、地鐵裡都行。
  \item \textbf{隱私更安全}：你的對話和資料不會上傳到雲端，完全在手機本地處理。
  \item \textbf{響應更快}：沒有網路延遲，AI回覆幾乎是即時的。
  \item \textbf{當前限制}：端側模型的能力還比不上雲端大模型，複雜任務仍需聯網。但隨著晶片AI算力的提升，端側模型會越來越強。
\end{itemize}

\booksubsection{系統級AI整合：}

\begin{itemize}
  \item \textbf{蘋果Apple Intelligence}：iOS 18引入的AI系統，整合寫作工具、影象魔法棒、Siri增強、通知摘要等功能。端側處理優先，複雜任務才上雲端。
  \item \textbf{華為盤古大模型}：整合到鴻蒙系統中，小藝助手獲得大模型能力，支援檔案摘要、智慧識屏、AI修圖等。
  \item \textbf{小米澎湃AI}：系統級AI能力，支援AI搜圖、AI摘要、AI實時字幕、端側大模型對話等。
  \item \textbf{三星Galaxy AI}：通話實時翻譯、筆記AI摘要、圖片生成編輯、搜尋圈選等功能。
\end{itemize}

\booksection{五、幾個典型應用場景與實操}

為了讓你更直觀地感受手機AI的“超能力”，下面我們用五個具體場景，展示AI是如何在日常生活中“出手相救”的：

\booksubsection{場景 1：實時翻譯——出國旅行的隨身翻譯官}

\textbf{實際表現}：在巴黎點餐不會法語？iPhone的“翻譯”App支援面對面雙人對話模式，雙方各看自己語言的字幕，甚至可以同聲傳譯。安卓陣營的三星Galaxy AI通話實時翻譯能讓你在電話裡直接說中文，對方聽到的是英文，反之亦然。

\textbf{背後原理}：語音識別（ASR） 機器翻譯（NMT） 語音合成（TTS）三步流水線，端到端延遲已可控制在1秒內。

\booksubsection{場景 2：AI消除——拯救被路人“毀掉”的遊客照}

\textbf{實際表現}：在景點拍了一張絕美的自拍，但背景裡擠滿了其他遊客？用“魔法消除”圈選路人，手機會自動識別並把路人“塗掉”，再智慧填充背景，效果常常令人驚豔。

\textbf{背後原理}：影象分割（Inpainting）模型先識別並刪除選中區域，然後根據周圍畫素和語義理解“腦補”出合理的填充內容。這背後的AI模型訓練資料量高達數百萬張圖片。

\booksubsection{場景 3：AI修圖——一鍵變“攝影大師”}

\textbf{實際表現}：拍了一張陰天的灰濛濛風景照？AI“魔法天空”可以一鍵把天空換成晴空、晚霞甚至星空。AI擴圖可以讓豎構圖變成橫構圖，手機“腦補”出畫面外的場景。

\textbf{背後原理}：基於生成式AI（Generative AI）的擴散模型（Diffusion Model），理解影象的語義內容後生成新畫素。這是當前AI影象處理的“最前沿”。

\booksubsection{場景 4：檔案處理——辦公族的效率倍增器}

\textbf{實際表現}：會議結束後，把會議錄音丟給AI，它能自動整理出結構化的會議紀要、提取待辦事項、生成摘要。一份10頁的英文PDF報告，AI可以在10秒內給你提煉出核心結論。

\textbf{背後原理}：長上下文大模型（Long Context LLM）可以“讀完”數萬字的檔案，進行多輪推理和資訊抽取。Apple Intelligence、華為小藝、WPS AI、騰訊元寶等都具備此能力。

\booksubsection{場景 5：AI搜圖與識物——你的“視覺搜尋引擎”}

\textbf{實際表現}：在街上看到一種不認識的花，相機一開，AI立刻告訴你“這是木槿花”。相簿裡翻找去年在海邊拍的照片，搜尋“大海”，AI自動把所有含海的照片找出來，不用手動翻。

\textbf{背後原理}：計算機視覺（CV）模型對圖片內容進行向量化（Embedding），再與文字查詢進行相似度匹配。這種“多模態搜尋”是手機相簿管理的一大飛躍。

\booksection{六、提示詞最佳化：讓 AI 聽你指揮的“溝通公式”}

AI 的回答質量，很大程度上取決於你如何“提問”。同樣一個問題，問法不同，AI 給的答案天差地別。下面我們介紹一套\textbf{最實用、最容易上手的“角色-背景-任務-限制”溝通思路}，並配以大量例項。

\begin{tcbinfo}{為什麼需要框架？}
很多人用 AI 像用百度一樣——敲幾個關鍵字進去。但 AI 不是搜尋引擎，它是一個“閱讀理解後回答”的語言模型。\textbf{給它更多上下文、說清你要什麼、限制它別跑偏}，它就能給你精準的結果。框架就是把你腦子裡的“好答案長什麼樣”翻譯成 AI 能理解的指令。

\end{tcbinfo}

\booksubsection{角色（Role）：告訴 AI 它是誰}

給 AI 一個具體的“身份”，它就會呼叫對應領域的專業知識和表達方式。\textbf{角色設定得越具體，輸出越專業}。

\booksubsection{實用示例：}

\begin{itemize}
  \item "你是一位\textbf{有 10 年經驗的小米手機維修工程師}。"（技術問題）
  \item "你是一位\textbf{擅長把複雜概念講簡單的中學物理老師}。"（科普問題）
  \item "你是一位\textbf{審美好、讀得多、有自己風格的圖書編輯}。"（寫作問題）
  \item "你是一位\textbf{每天給上班族做早餐的烹飪博主}。"（生活問題）
  \item "你是一位\textbf{耐心、溫和、永遠不嘲笑新手的中醫}。"（健康問題）
\end{itemize}

\textbf{為什麼有效？}角色設定了 AI 的\textbf{知識範圍}、\textbf{表達風格}和\textbf{專業深度}。沒有角色的 AI 是個“全才但啥都不精”的通才；有了角色，它能“進入”那個身份，像那位專業人士一樣思考。

\booksubsection{背景（Background）：告訴 AI 你的情況}

把“你是誰、你的現狀、你的約束條件”提前說清楚，AI 就能給出\textbf{真正適合你}的答案，而不是泛泛而談。

\booksubsection{實用示例：}

\begin{itemize}
  \item "我的手機是\textbf{小米 14，買了半年，最近玩遊戲時發燙嚴重}。"（裝置情況）
  \item "我是\textbf{一個 60 歲的退休阿姨，眼睛不太好，不太會用智慧手機}。"（使用者畫像）
  \item "我正在\textbf{準備 2026 年 6 月的英語六級考試，基礎一般}。"（學習背景）
  \item "我開了一家\textbf{三線城市的小咖啡店，月流水 3 萬，僱了兩個員工}。"（商業背景）
  \item "我\textbf{剛生完寶寶在家坐月子，睡眠嚴重不足}。"（生活階段）
\end{itemize}

\textbf{為什麼有效？}背景資訊讓 AI 知道\textbf{“為誰服務”}和\textbf{“在什麼條件下”}。同一道題，給學生和給家長講法不同；給城裡人和給鄉鎮人建議也不同。背景缺失，AI 只能給“最大公約數”的答案，往往不痛不癢。

\booksubsection{任務（Task）：明確告訴 AI 做什麼}

這一步要用\textbf{動詞開頭}（寫、列、對比、翻譯、解釋、總結……），把任務描述得\textbf{具體、可執行}。含糊的任務只會換來含糊的答案。

\booksubsection{實用示例：}

\begin{itemize}
  \item "請\textbf{列出 5 個我應該立刻關閉的權限}，並說明原因。"（動作明確）
  \item "請\textbf{把下面這段英文郵件翻譯成中文}，保持商務口吻。"（任務具體）
  \item "請\textbf{對比這四款 3000 元價位的手機}，給出一張優缺點對比表。"（格式明確）
  \item "請\textbf{用小學生能懂的話解釋'什麼是 5G'}。"（風格明確）
  \item "請\textbf{給我的早餐店寫三條短影片文案}，分別針對上班族、情侶、帶孩子的家長。"（細分明確）
\end{itemize}

\textbf{為什麼有效？}任務部分是 AI 行動的“指令中樞”。\textbf{動詞越清晰、要求越具體，AI 越知道“交付什麼”}。模糊的“弄一下”、“搞個”、“看看”是 AI 大忌——它會揣測你的意思，結果常常跑偏。

\booksubsection{限制（Constraints）：告訴 AI 不要做什麼}

“限制”聽起來像“約束”，其實是\textbf{讓 AI 收斂到最適合你的輸出}的利器。常見維度有：字數、格式、語氣、避免項、必含項。

\booksubsection{實用示例：}

\begin{itemize}
  \item "\textbf{不超過 200 字，分點列出}。"（長度+格式）
  \item "\textbf{語氣要親切自然，像朋友聊天}，不要用'首先、其次、再次'這種官話。"（語氣）
  \item "\textbf{不要使用專業術語}，遇到專業名詞要順便解釋。"（用詞限制）
  \item "\textbf{每段不超過 3 行}，方便手機閱讀。"（排版）
  \item "\textbf{必須包含：使用場景、優缺點、我的預算匹配度}。"（必含項）
  \item "\textbf{不要給具體投資建議}，只解釋原理。"（邊界）
\end{itemize}

\textbf{為什麼有效？}沒有限制，AI 容易“自由發揮”——要麼太學術、要麼太長、要麼跑題。明確的限制=明確的“交付標準”，雙方對齊預期，結果才可控。

\booksubsection{提示詞質量評估：好提示詞的 5 個特徵}

寫好提示詞不是一次性的事，\textbf{寫完先自查這五條}，再發給 AI：

\begin{enumerate}
  \item \textbf{具體而非抽象}：不是"寫篇文章"，而是"寫一篇 800 字、3 個小節、面向新手媽媽的手機選購指南"。
  \item \textbf{包含背景而非孤立}：告訴 AI"我是誰、為什麼問這個"，它才能給出"為你定製"的答案。
  \item \textbf{有明確輸出格式}：是表格、列表、段落、還是 JSON？提前說。
  \item \textbf{限制到位}：字數、風格、避項，三選一就有明顯提升。
  \item \textbf{可被 AI 驗證}：問完之後，讓 AI 複述一遍你的要求；說不清楚就是沒說清楚。
\end{enumerate}

\booksubsection{常見問題與解法}

\begin{center}
\begin{tabularx}{\linewidth}{|X|X|}
\hline
\textbf{症狀} & \textbf{解法} \tabularnewline
\hline
回答太泛、不實用 & 加背景（你是誰/什麼場景） \tabularnewline
回答太學術/太官方 & 加語氣限制 + 拒絕"官話"清單 \tabularnewline
回答跑題/答非所問 & 把任務用動詞重新說一遍，並要求"只回答這個" \tabularnewline
回答太長看不完 & 硬性限制字數（200 字內） \tabularnewline
AI 一本正經胡說 & 加"如不確定請標註'我不確定'" + 重要事實自查 \tabularnewline
第一版不滿意 & 不要重寫，在原提示詞後追加"再…一些" \tabularnewline
\hline
\end{tabularx}
\end{center}

\booksubsection{三個完整場景案例}

\begin{tcbinfo}{場景一：生活助手（早餐搭配建議）}
\textbf{提示詞}：

你是一位擅長營養搭配的早餐博主（\textbf{角色}）。我是一位上班族媽媽，每天早上只有 15 分鐘做早餐，家裡有 6 歲的孩子和 60 歲的婆婆（\textbf{背景}）。請\textbf{列出 7 天不重樣的早餐搭配}（\textbf{任務}）。每份早餐要求：(1) 準備時間不超過 15 分鐘 (2) 營養均衡（蛋白+碳水+蔬果）(3) 老少皆宜 (4) 食材在普通超市能買到（\textbf{限制}）。

\textbf{為什麼有效}：角色定專業度，背景給三類人共性，任務明確交付物，4 條限制確保可執行。

\end{tcbinfo}

\begin{tcbinfo}{場景二：學習輔助（英語六級備考）}
\textbf{提示詞}：

你是一位\textbf{經驗豐富的英語六級老師}（角色）。我是一名普通本科生，\textbf{四級 480 分，六級目標 500，還有 3 個月}（背景）。請\textbf{為我制定一份可執行的 12 周備考計劃}（任務）。要求：(1) 每週明確聽說讀寫的時間分配 (2) 給出每個月的重點突破方向 (3) 推薦 3 本合適的複習資料（中文說明理由） (4) 表格形式輸出（限制）。

\textbf{為什麼有效}：角色定專業背景，背景給出分數和時間錨點，任務明確“分階段計劃”，限制要求“表格化”和“可執行”。

\end{tcbinfo}

\begin{tcbinfo}{場景三：工作效率（會議紀要整理）}
\textbf{提示詞}：

你是一位\textbf{資深行政助理}（角色）。我剛開完一個 1 小時的部門會議，\textbf{討論 Q2 營銷推廣方案}（背景）。\textbf{請把以下會議記錄整理成結構化紀要}（任務）。要求：(1) 分“決議事項 / 待辦事項 / 風險與異議”三部分 (2) 待辦事項標出負責人和截止時間（如未提到則寫“待確認”） (3) 總字數控制在 600 字內 (4) 用 markdown 格式輸出（限制）。

\textbf{為什麼有效}：角色定行政視角，背景給會議主題，任務明確結構化輸出，4 條限制讓結果直接可複用。

\end{tcbinfo}

\textbf{進階玩法：讓 AI 成為你的“私人助理”}\par

試試這些進階用法，體驗 AI 的真正實力：

\begin{itemize}
  \item \textbf{多輪對話最佳化}：對 AI 的回答不滿意時，可以說"這個太專業了，請用小學生能懂的話再解釋一次"。
  \item \textbf{鏈式思考（CoT）}：讓 AI"一步步思考"，回答質量通常更高："請一步步分析這道數學題"。
  \item \textbf{多 AI 協作}：用 A 寫初稿，B 潤色，C 翻譯——組合使用效果更佳。
  \item \textbf{系統提示詞}：在 AI 設定中預設你的身份、偏好、忌諱等，讓回答更貼合你的需求。
\end{itemize}

\booksection{七、AI的侷限與挑戰：保持清醒的判斷力}

AI雖然強大，但絕非“無所不能”。清醒地認識它的侷限，才能更好地使用它：

\begin{itemize}
  \item \textbf{"幻覺"問題}：AI可能一本正經地胡說八道（專業術語叫Hallucination）。它會編造不存在的歷史事件、錯誤的引言、虛假的統計資料。重要資訊必須人工核實。
  \item \textbf{知識截止}：大多數AI的訓練資料有截止日期（如2024年初），它不知道之後發生的事。
  \item \textbf{推理能力有限}：複雜的多步推理、數學證明、邏輯陷阱，AI仍然容易出錯。
  \item \textbf{隱私風險}：在雲端AI服務中，你的對話可能被用於模型訓練。敏感資訊應優先使用端側AI。
  \item \textbf{偏見與價值觀}：AI可能帶有訓練資料中的偏見，需要保持獨立思考。
\end{itemize}

\begin{tcbinfo}{AI使用小貼士}
\begin{itemize}
  \item AI生成的內容可能不準確（"幻覺"問題），重要資訊務必核實。
  \item 不要向AI輸入個人隱私資訊（身份證號、密碼、銀行資訊等）。
  \item 善用提示詞（Prompt）：越具體、越清晰的描述，AI的回答質量越高。
  \item AI是工具不是權威，保持獨立思考和判斷力。
  \item 遇到關鍵決策（醫療、法律、財務），諮詢專業人士而非依賴AI。
\end{itemize}

\end{tcbinfo}

\begin{tcbpractice}{來挑戰：三道AI實戰題}
\begin{enumerate}
  \item \textbf{挑戰一：複雜指令}：對語音助手說："明早7點提醒我帶傘，並把今天下午3點的日曆改到4點。" 看看它能否一次完成多步操作。
  \item \textbf{挑戰二：AI修圖}：拍一張有路人的遊客照，使用"魔法消除"或"AI橡皮擦"功能，把路人清除乾淨。觀察AI對背景的"腦補"效果。
  \item \textbf{挑戰三：提示詞最佳化}：先用簡單提示詞"寫一段旅遊攻略"問AI，再用CRISPE框架寫一個詳細提示詞。比較兩次回答的質量差異。
\end{enumerate}

\end{tcbpractice}

\begin{tcbinfo}{案例 2：換機資料遷移完整流程}
\begin{itemize}
  \item \textbf{場景}：iPhone 13 小米 15 跨品牌換機。
  \item \textbf{步驟}：(1) 舊機 iCloud/微信自帶備份 (2) 退出 Apple ID + 關閉"查詢" (3) 新機用"小米換機"App 掃碼 (4) 選擇遷移項 (5) 微信單獨用自帶"聊天記錄遷移" (6) 驗證關鍵資料（通訊錄/相簿/微信） (7) 舊機加密擦除。
  \item \textbf{預期結果}：資料零丟失，舊機可安全轉賣。
\end{itemize}

\end{tcbinfo}
